% GENERATED FILE. Edit canonical lessons, then run npm run generate:books.
% canonical-source-hash: fnv1a64:c43665d7e3b2218e
% canonical-lessons: PA-C196-work-command-return, PA-C196-shuru, PA-C196-shuru-review, PA-C196-start-work-frame, PA-C196-payoff

\chapter{Start the Work}
\label{ch:pa-a2-shuru}
\hlchaptermodality{196}

\begin{chapteropening}
\textbf{By the end of this chapter:} \emph{I can recognise and use the familiar word for start in one bounded polite work instruction.}

Recognise and produce the start word in one bounded polite work instruction through listening, speaking, reading, and writing.
\end{chapteropening}

\section[\texorpdfstring{\pa{ਕੰਮ} · \pa{ਕਰੋ} (kamm; karo)}{ਕੰਮ · ਕਰੋ (kamm; karo)}]{Return to work and a polite command}

\label{lesson:PA-C196-work-command-return}



\begin{quote}
\emph{Hear:} \emph{kamm; karo}; YOU SAY: \textbf{\pa{ਕੰਮ}; \pa{ਕਰੋ}}
\end{quote}

\begin{grammarlens}[title={Grammar Lens}]
Keep the familiar work word first and the familiar polite command last. The middle slot stays empty until the next lesson owns its new word.
\end{grammarlens}

\subsection*{Your turn}
\begin{itemize}
\raggedright
  \item \emph{Hear:} \emph{kamm}; YOU SAY: \textbf{\pa{ਕੰਮ}}
  \item \emph{Hear:} \emph{karo}; YOU SAY: \textbf{\pa{ਕਰੋ}}
  \item Say the first word, pause for the empty middle, then say the last.
\end{itemize}

\begin{tcolorbox}[breakable,colback=teal!4,colframe=teal!35!black,title={Before you move on}]
Produce the two familiar words around the empty middle slot. This opener owns no new word.
\end{tcolorbox}
\section[\texorpdfstring{\pa{ਸ਼ੁਰੂ} (shuruu)}{ਸ਼ੁਰੂ (shuruu)}]{One noun for a beginning or start}

\label{lesson:PA-C196-shuru}



\begin{quote}
Retrieve the familiar first and last words of the work instruction.
\end{quote}

\begin{grammarlens}[title={Grammar Lens: direct owner}]
\begin{quote}
\textbf{\pa{ਸ਼ੁਰੂ}} — \emph{shurū} — beginning; start
\end{quote}

Hear and say the whole word before reading it. Keep it as a whole word rather than inventing smaller Punjabi pieces. This lesson owns the word alone; the complete instruction comes only after a separate retrieval.
\end{grammarlens}

\begin{tcolorbox}[breakable,skin=enhanced,colback=violet!6,colframe=violet!45!black,boxrule=0.5pt,arc=1mm,left=6pt,right=6pt,top=4pt,bottom=4pt,fonttitle=\bfseries,title={Across the family}]
The Punjabi University dictionary records this Punjabi noun for a beginning or start; its index uses the equivalent precomposed letter spelling. Wiktionary records the Punjabi word as borrowed from Arabic. Keep that source history separate from present-day Punjabi use, and learn the Punjabi form as one whole. Sources: \href{https://dic.learnpunjabi.org/default.aspx?look=\%E0\%A8\%B6\%E0\%A9\%81\%E0\%A8\%B0\%E0\%A9\%82}{Punjabi University Punjabi-English Dictionary}; \href{https://en.wiktionary.org/wiki/\%E0\%A8\%B8\%E0\%A8\%BC\%E0\%A9\%81\%E0\%A8\%B0\%E0\%A9\%82}{Wiktionary: \pa{ਸ਼ੁਰੂ}}.
\end{tcolorbox}

\subsection*{Your turn}
\begin{itemize}
\raggedright
  \item \emph{Hear:} \emph{shurū}; YOU SAY: \textbf{\pa{ਸ਼ੁਰੂ}}
  \item \emph{Read it:} \textbf{\pa{ਸ਼ੁਰੂ}}
  \item \emph{Write it:} copy \textbf{\pa{ਸ਼ੁਰੂ}} once, checking the whole word
\end{itemize}

\begin{tcolorbox}[breakable,colback=teal!4,colframe=teal!35!black,title={Before you move on}]
Produce the noun for “beginning” or “start” without looking.
\end{tcolorbox}
\section[\texorpdfstring{\pa{ਸ਼ੁਰੂ} (shuruu)}{ਸ਼ੁਰੂ (shuruu)}]{Retrieve the new word on its own}

\label{lesson:PA-C196-shuru-review}



\begin{quote}
\emph{Hear:} \emph{shurū}; YOU SAY: \textbf{\pa{ਸ਼ੁਰੂ}}
\end{quote}

\subsection*{Your turn}
\begin{itemize}
\raggedright
  \item \emph{Hear:} \emph{beginning; start}; YOU SAY: \textbf{\pa{ਸ਼ੁਰੂ}}
  \item \emph{Read it:} \textbf{\pa{ਸ਼ੁਰੂ}}
  \item \emph{Write it:} cover and write \textbf{\pa{ਸ਼ੁਰੂ}}
\end{itemize}

\begin{grammarlens}[title={Grammar Lens}]
Keep one noun and one meaning here. Broader location and verb patterns stay outside this bounded lesson.
\end{grammarlens}

\begin{tcolorbox}[breakable,colback=teal!4,colframe=teal!35!black,title={Before you move on}]
Say and then write the word once from memory.
\end{tcolorbox}
\section[\texorpdfstring{\pa{ਕੰਮ} \pa{ਸ਼ੁਰੂ} \pa{ਕਰੋ} (kamm shuruu karo)}{ਕੰਮ ਸ਼ੁਰੂ ਕਰੋ (kamm shuruu karo)}]{Put the new word into one work instruction}

\label{lesson:PA-C196-start-work-frame}



\begin{quote}
Retrieve \textbf{\pa{ਕੰਮ}}, then \textbf{\pa{ਸ਼ੁਰੂ}}, then \textbf{\pa{ਕਰੋ}}.
\end{quote}

\begin{grammarlens}[title={Grammar Lens}]
\begin{quote}
\textbf{\pa{ਕੰਮ} \pa{ਸ਼ੁਰੂ} \pa{ਕਰੋ}} — \emph{kamm shurū karo} — start the work
\end{quote}

Keep the new noun whole between familiar \textbf{\pa{ਕੰਮ}} and \textbf{\pa{ਕਰੋ}}. Familiar \textbf{\pa{ਕਰੋ}} carries the polite command in the already-owned noun-plus-\textbf{\pa{ਕਰਨਾ}} pattern. Gill and Gleason describe such noun/adjective-plus-verb conjuncts as conventional combinations: this chapter teaches only this one instruction, not every possible combination, tense, or function. Source: Harjeet Singh Gill and Henry A. Gleason, Jr., \emph{A Reference Grammar of Punjabi} (1969), pp. 122 and 124, \href{https://pt.learnpunjabi.org/assets/A\%20Reference\%20Grammar\_Final.pdf}{conjunct verbs}.
\end{grammarlens}

\subsection*{Your turn}
\begin{itemize}
\raggedright
  \item \emph{Hear:} \emph{kamm shurū karo}; YOU SAY: \textbf{\pa{ਕੰਮ} \pa{ਸ਼ੁਰੂ} \pa{ਕਰੋ}}
  \item \emph{Read it:} \textbf{\pa{ਕੰਮ} \pa{ਸ਼ੁਰੂ} \pa{ਕਰੋ}}
  \item \emph{Write it:} cover and write \textbf{\pa{ਕੰਮ} \pa{ਸ਼ੁਰੂ} \pa{ਕਰੋ}}
\end{itemize}

\begin{tcolorbox}[breakable,colback=teal!4,colframe=teal!35!black,title={Before you move on}]
Produce the complete polite instruction.
\end{tcolorbox}
\section[\texorpdfstring{\pa{ਕੰਮ} \pa{ਸ਼ੁਰੂ} \pa{ਕਰੋ} (kamm shuruu karo)}{ਕੰਮ ਸ਼ੁਰੂ ਕਰੋ (kamm shuruu karo)}]{A bounded start-word payoff}

\label{lesson:PA-C196-payoff}



\begin{quote}
Retrieve the new word alone before rebuilding its instruction.
\end{quote}

\subsection*{Your turn: four skills}
\begin{itemize}
\raggedright
  \item \emph{Hear:} \emph{kamm shurū karo}; YOU SAY: \textbf{\pa{ਕੰਮ} \pa{ਸ਼ੁਰੂ} \pa{ਕਰੋ}}
  \item \emph{Read it:} \textbf{\pa{ਕੰਮ} \pa{ਸ਼ੁਰੂ} \pa{ਕਰੋ}}
  \item \emph{Write it:} cover the line and write \textbf{\pa{ਕੰਮ} \pa{ਸ਼ੁਰੂ} \pa{ਕਰੋ}}
\end{itemize}

\begin{tcolorbox}[breakable,colback=teal!4,colframe=teal!35!black,title={Before you move on}]
Say the new word, then produce the complete work instruction.

This is the nineteenth prerequisite tranche exposed by Punjabi A2 mock drafting. It does not unblock Punjabi A2 mock 1 by itself. It is not an A2 mock, not A2 readiness or attainment, and not external or human validation.
\end{tcolorbox}
